%% =====================================================================
%%  T-02-03  Patience as Position
%%  -------------------------------------------------------------------
%%  One idea per file. Core is mandatory. Slots in canonical order.
%%  Max 700 words. No layout commands. No filler. New source material
%%  for this entry goes in sources/<BOOK>/T-02-03.tex -- never by
%%  rewriting this file (docs/RULES.md R0.1).
%% =====================================================================
%% @kind: topic
%% @id: T-02-03
%% @title: Patience as Position
%% @subtitle: Hurry is a tell; waiting is a lever
%% @chapter: c02
%% @order: 30
%% @status: stable
%% @sources: B3
%% @updated: 2026-09-19

\topic{T-02-03}{Patience as Position}{Hurry is a tell; waiting is a lever}

\begin{core}
Apparent haste signals that you need the outcome more than the other person does, and that signal re-prices the whole exchange against you. Patience is therefore not a virtue here but a position: it keeps your need invisible and keeps the counterpart's need on the table.
\end{core}
\begin{mechanism}
Two mechanisms run together. Perceptually, hurry communicates dependence, and dependence is what the whole of Part III is about concealing. Practically, most situations have a moment at which they become cheap to move --- before it, effort is wasted; after it, the opportunity has closed. Detecting that moment requires having watched for longer than the other party expected you to. Impatience collapses both: it reveals the need and it spends the effort at the wrong time.
\end{mechanism}
\begin{tells}
  \item You have raised your own offer before being asked to.
  \item You are checking on something far more often than its importance warrants.
  \item You feel that a delay is a loss rather than a neutral event.
  \item The other party has begun setting the pace of the exchange.
\end{tells}
\begin{moves}
  \item Set a horizon before you start, and treat everything inside it as available time rather than lost time.
  \item Match the counterpart's response interval instead of beating it.
  \item Use the waiting period to gather information, which is the only productive use of it.
  \item Act decisively at the moment the situation turns, and not before --- standing back while the time is unripe and moving hard when it is ripe are the same skill.
\end{moves}
\begin{counters}
  \item Impose a deadline and see whether it holds. Manufactured urgency collapses under a request to wait; real urgency does not.
  \item Refuse to let pace be set by the other party's calendar. Ask for the reason behind the date.
  \item If you are being worn down by delay rather than by argument, name it: a tactic that depends on your fatigue stops working when it is identified.
\end{counters}
\begin{cost}
Patience requires that you can afford to wait, and both sources note that the person with high fixed costs cannot. Waiting is a lever only for the party with the smaller exposure, which means the honest first move is often to reduce your exposure before entering the exchange rather than to sit tight inside it. Patience also decays into drift if there is no defined trigger for acting.
\end{cost}

\sources{B3}
\seesources
\seealso{T-02-01, T-11-01, T-17-02}
